\title{Visual Prompt Tuning}

\begin{abstract}
The current \textit{modus operandi} in adapting pre-trained models involves updating all the backbone parameters, i.e. , full fine-tuning. This paper introduces Visual Prompt Tuning (VPT) as an efficient and effective alternative to full fine-tuning for large-scale Transformer models in vision. Taking inspiration from recent advances in efficiently tuning large language models, VPT introduces only a small amount (less than 1\% of model parameters) of trainable parameters in the input space while keeping the model backbone frozen. Via extensive experiments on a wide variety of downstream recognition tasks, we show that VPT achieves significant performance gains compared to other parameter efficient tuning protocols. Most importantly, VPT even outperforms full fine-tuning in many cases across model capacities and training data scales, while reducing per-task storage cost. Code is available at \href{https://github.com/kmnp/vpt}{\texttt{github.com/kmnp/vpt}}.
\end{abstract}

\section{Introduction}\label{sec: intro_luming}

For a variety of recognition applications, the most accurate results are now obtained by adapting large \emph{foundation models} pre-trained on massive curated or raw data, a finding that mirrors developments in natural language processing (NLP)~\cite{bommasani2021opportunities}.\footnote{As pointed out in~\cite{bommasani2021opportunities}, all state-of-the-art models in contemporary NLP are now powered by a few Transformer-based models (e.g. , BERT~\cite{devlin-etal-2019-bert}, T5~\cite{2020t5}, BART~\cite{lewis2020bart}, GPT-3~\cite{brown2020gpt3}) This also applies to vision-language field recently, i.e. , CLIP~\cite{radford2021learning}.} At first glance,this is a success story: one can make rapid progress on multiple recognition problems simply by leveraging the latest and greatest foundation model. In practice, however, \emph{adapting} these large models to downstream tasks presents its own challenges. The most obvious (and often the most effective) adaptation strategy is \emph{full fine-tuning} of the pre-trained model on the task at hand, end-to-end. However, this strategy requires one to store and deploy a separate copy of the backbone parameters for every single task. This is an expensive and often infeasible proposition, especially for modern \emph{Transformer}-based architectures, which are significantly larger than their convolutional neural networks (ConvNet) counterparts, e.g. , ViT-Huge~\cite{dosovitskiy2020vit} (632M parameters) \emph{vs}\onedot~ResNet-50~\cite{he2016rn} (25M parameters). We therefore ask, \textbf{what is the best way to adapt large pre-trained Transformers to downstream tasks in terms of effectiveness and efficiency}?

One straightforward approach is to turn to other strategies that we have perfected for adapting ConvNets to new tasks, as in~\ref{fig:teaser}(a).
A popular approach is to fine-tune only a subset of the parameters, such as the classifier head~\cite{wslimageseccv2018,jia2021exploring,chen2021mocov3} or the bias terms~\cite{cai2020tinytl}. Prior research has also looked at adding additional residual blocks  (or \emph{adapters}) to the backbone~\cite{rebuffi2018efficient,zhang2020side}. One could implement similar strategies for Transformers. However, in general these strategies \emph{under-perform} full fine-tuning in accuracy.

We explore a different route in this paper. Instead of altering or fine-tuning the pre-trained Transformer itself, we modify the \emph{input} to the Transformer. Drawing inspiration from the recent advances on Prompting in NLP~\cite{liu2021pre,li-liang-2021-prefix,lester-etal-2021-power,liu2021p}, we propose a new simple and efficient method to adapt transformer models for downstream vision tasks (\ref{fig:teaser}(b)), namely \textbf{Visual-Prompt Tuning} (VPT). Our method only introduces a small amount of task-specific learnable parameters into the input space while freezing the entire pre-trained Transformer backbone during downstream training. In practice, these additional parameters are simply prepended into the input sequence of each Transformer layer and learned together with a linear head during fine-tuning.

On 24 downstream recognition tasks spanning different domains using a pre-trained ViT backbone, VPT beats all other transfer learning baselines, even surpassing full fine-tuning in 20 cases, while maintaining the advantage of storing remarkably fewer parameters (less than 1\% of backbone parameters) for each individual task (\ref{fig:teaser}(c)). 
This result demonstrates the distinctive strength of \emph{visual} prompting: whereas in NLP, prompt tuning is only able to \emph{match} full fine-tuning performance under certain circumstances~\cite{lester-etal-2021-power}. VPT is especially effective in the low-data regime, and maintains its advantage across data scales. Finally, VPT is competitive for a range of Transformer scales and designs (ViT-Base/Large/Huge, Swin). Put together, our results suggest that VPT is one of the most effective ways of adapting ever-growing vision backbones.

\section{Related Work}\label{sec:related}

\textbf{Transformer}
models~\cite{vaswani2017attention} have gained huge success in NLP~\cite{devlin-etal-2019-bert,2020t5,brown2020gpt3}. The triumph of the Transformer architecture also extends to various computer vision tasks, including image classification~\cite{dosovitskiy2020vit,liu2021swin}, object detection~\cite{carion2020end,li2021benchmarking}, semantic and panoptic segmentation~\cite{strudel2021segmenter,zheng2020rethinking,wang2021max}, video understanding~\cite{girdhar2019video,wang2022bevt,feichtenhofer2022masked} and few-shot learning~\cite{doersch2020crosstransformers}, surpassing previous state-of-the-art approaches. Transformers are also being widely used in recent self-supervised pre-training methods~\cite{chen2021mocov3,he2021mae,bao2021beit}. Given their superior performance and much larger scale compared to ConvNets, how to efficiently adapt Transformers to different vision tasks remains an important open problem. Our proposed VPT provides a promising path forward.

\textbf{Transfer learning} has been extensively studied for vision tasks in the context of ConvNets~\cite{zhuang2020comprehensive} and many techniques have been introduced including side tuning~\cite{zhang2020side}, residual adapter~\cite{rebuffi2017learning}, bias tuning~\cite{cai2020tinytl}, etc. . Relatively little attention has been paid to vision Transformers adaptation and how well these aforementioned methods perform on this brand new type of architecture remains unknown. On the other hand, given the dominance of large-scale pre-trained Transformer-based Language Models (LM)~\cite{devlin-etal-2019-bert,2020t5,brown2020gpt3}, many approaches~\cite{he2022towards,guo2020parameter,hu2021lora} have been proposed to efficiently fine-tune LM for different downstream NLP tasks~\cite{wang2018glue,wang2019superglue}. Among them, we focus on the following two representative methods in our experiments for benchmarking purposes: Adapters~\cite{pfeiffer2020AdapterHub} and BitFit~\cite{zaken2021bitfit}.

Adapters~\cite{houlsby2019parameter} insert extra lightweight modules inside each Transformer layer. One adapter module generally consists of a linear down-projection, followed by a nonlinear activation function, and a linear up-projection, together with a residual connection~\cite{pfeiffer2020adapterfusion,pfeiffer2020AdapterHub}. Instead of inserting new modules, \cite{cai2020tinytl} proposed to update the bias term and freeze the rest of backbone parameters when fine-tuning ConvNets. BitFit~\cite{bao2021beit} applied this technique to Transformers and verified its effectiveness on LM tuning. Our study demonstrates that VPT, in general, provides improved performance in adapting Transformer models for vision tasks, relative to the aforementioned two well-established methods in NLP.

\textbf{Prompting}~\cite{liu2021pre} originally refers to prepending language instruction to the input text so that a pre-trained LM can ``understand'' the task. With manually chosen prompts, GPT-3 shows strong generalization to downstream transfer learning tasks even in the few-shot or zero-shot settings~\cite{brown2020gpt3}. In addition to the follow-up works on how to construct better prompting texts~\cite{shin2020autoprompt,jiang2020can}, recent works propose to treat the prompts as task-specific continuous vectors and directly optimize them via gradients during fine-tuning, namely Prompt Tuning~\cite{li-liang-2021-prefix,lester-etal-2021-power,liu2021p}. Compared to full fine-tuning, it achieves comparable performance but with 1000$\times$ less parameter storage. Although prompting has also been applied to vision-language models recently~\cite{radford2021learning,zhou2021learning,ju2021prompting,yao2021cpt,ge2022domain}, prompting is still limited to the input of \emph{text} encoders. Due to the disparity between vision and language modalities, in this paper we ask: can the same method can be applied successfully to image encoders? We are the first work (see related concurrent works~\cite{sandler2022fine,wang2022learning,conder2022efficient,bahng2022visual}) to tackle this question and investigate the generality and feasibility of visual prompting via  \emph{extensive} experiments spanning multiple kinds of recognition tasks across multiple domains and backbone architectures.

\section{Approach}\label{sec:method}

We propose Visual-Prompt Tuning (\textsc{VPT}) for adapting large pre-trained vision Transformer models. 
\textsc{VPT} injects a small number of learnable parameters into Transformer's input space and keeps the backbone frozen during the downstream training stage. The overall framework is presented in \ref{fig:method}.
We first define the notations in~\ref{subsec:method_pre}, then describe VPT formally in~\ref{subsec:method_vp}.

\subsection{Preliminaries}
\label{subsec:method_pre}

For a plain Vision Transformer (ViT)~\cite{dosovitskiy2020vit} with $N$ layers, an input image is divided into $m$ fixed-sized patches 
$\{I_j\in\mathbb{R}^{3\times h\times w}\mid j\in\mathbb{N}, 1\le j\le m\}$. $h, w$ are the height and width of the image patches. Each patch is then first embedded into $d$-dimensional latent space with positional encoding:

\begin{align}
\label{eq:embed}
\vec{e}_0^j = \texttt{Embed}(I_j) &&\vec{e}_0^j\in\mathbb{R}^{d}, j = 1,2, \ldots m \;\;.
\end{align}

We denote the collection of image patch embeddings, 
$\vec{E}_{i}=\{\vec{e}_i^j\in\mathbb{R}^d\mid j\in\mathbb{N}, 1\le j\le m\}$,  
as inputs to the ($i$+$1$)-th Transformer layer ($L_{i+1}$). Together with an extra learnable classification token ($\texttt{[CLS]}$), the whole ViT is formulated as:

\begin{align}
\label{eq:vit}
    [\vec{x}_i, \vec{E}_i] &= L_i([\vec{x}_{i-1}, \vec{E}_{i-1}])  &&i=1, 2, \ldots, N \\
    \vec{y} &= \texttt{Head}(\vec{x}_N)\;\;,
\end{align}

where $\vec{x}_{i}\in\mathbb{R}^{d}$ denote $\texttt{[CLS]}$'s embedding at $L_{i+1}$'s input space. 
$[\cdot,\cdot]$ indicates stacking and concatenation on the sequence length dimension, i.e. , $[\vec{x}_{i}, \vec{E}_{i}]\in\mathbb{R}^{(1+m)\times d}$. Each layer $L_i$ consists of Multiheaded Self-Attention (MSA) and Feed-Forward Networks (FFN) together with LayerNorm~\cite{ba2016layer} and residual connections~\cite{he2016rn}. A neural classification head is used to map the final layer's $\texttt{[CLS]}$ embedding, $\vec{x}_{N}$, into a predicted class probability distribution $\vec{y}$.\footnote{Some Transformer architectures in Vision such as~Swin~\cite{liu2021swin} do not use $\texttt{[CLS]}$ and treat global pooled $\vec{E}_N$ as input for \texttt{Head}. We follow their designs when adapting VPT to these Transformer variants. See~\ref{supsec:detail} for more details.}

\subsection{Visual-Prompt Tuning (\textsc{VPT})}
\label{subsec:method_vp}

Given a pre-trained Transformer model, we introduce a set of $p$ continuous embeddings of dimension $d$, i.e. {},~\emph{prompts}, in the input space after the \texttt{Embed} layer. Only the task-specific prompts are being updated during fine-tuning, while the Transformer backbone is kept frozen. Depending on the number of Transformer layers involved, our approach has two variants, \textsc{VPT}-\textsc{shallow} and \textsc{VPT}-\textsc{deep}, as shown in~\ref{fig:method}.

\subsubsection{VPT-Shallow.}
Prompts are inserted into the first Transformer layer $L_1$ only. Each prompt token is a learnable $d$-dimensional vector. A collection of $p$ prompts is denoted as 
$\vec{P}=\{\vec{p}^k\in\mathbb{R}^d\mid k\in\mathbb{N}, 1\le k\le p\}$, the shallow-prompted ViT is:

\begin{align}
\label{eq:shallow}
    [\vec{x}_{1}, \vec{Z}_{1}, \vec{E}_{1}] &= \textcolor{prompt_blue}{L_1}([\textcolor{prompt_blue}{\vec{x}_{0}}, \textcolor{prompt_red}{\vec{P}},\vec{E}_{0}])   \\
    [\vec{x}_{i}, \vec{Z}_{i}, \vec{E}_{i}] &= \textcolor{prompt_blue}{L_i}([\vec{x}_{i-1}, \vec{Z}_{i-1},\vec{E}_{i-1}])  &&i=2, 3, \ldots, N\\
    \vec{y} &= \textcolor{prompt_red}{\texttt{Head}}(\vec{x}_N)\;\;,
\end{align}

where $\vec{Z}_{i}\in\mathbb{R}^{p\times d}$ represents the features computed by the $i$-th Transformer layer, and $[\vec{x}_{i}, \vec{Z}_{i}, \vec{E}_{i}]\in\mathbb{R}^{(1+p+m)\times d}$. The colors \textcolor{prompt_red}{$\bullet$} and \textcolor{prompt_blue}{$\bullet$} indicate \textcolor{prompt_red}{learnable} and \textcolor{prompt_blue}{frozen} parameters, respectively. Notably for ViT, $\vec{x}_N$ is invariant to the location of prompts since they are inserted after positional encoding, e.g. , $[\vec{x}_{0}, \vec{P},\vec{E}_{0}]$ and $[\vec{x}_{0}, \vec{E}_{0}, \vec{P}]$ are mathematically equivalent. This also applies to VPT-Deep.

\subsubsection{VPT-Deep.}
Prompts are introduced at \emph{every} Transformer layer's input space. For ($i$+$1$)-th Layer $L_{i+1}$, we denote the collection of input learnable prompts as 
$\vec{P}_{i}=\{\vec{p}_i^k\in\mathbb{R}^d\mid k\in\mathbb{N}, 1\le k\le m\}$. The deep-prompted ViT is formulated as:

\begin{align}
\label{eq:deep}
    [\vec{x}_i, \underline{}, \vec{E}_i] &= \textcolor{prompt_blue}{L_i}([\vec{x}_{i-1}, \textcolor{prompt_red}{\vec{P}_{i-1}},\vec{E}_{i-1}]) &&i=1, 2, \ldots, N\\
    \vec{y} &= \textcolor{prompt_red}{\texttt{Head}}(\vec{x}_N)\;\;.
\end{align}

\subsubsection{Storing Visual Prompts.} \textsc{VPT} is beneficial in presence of multiple downstream tasks. We only need to store the learned prompts and classification head for each task and re-use the original copy of the pre-trained Transformer model, significantly reducing the storage cost. For instance, given a ViT-Base with 86 million (M) parameters and $d=768$, 50 shallow prompts and deep prompts yield additional $p\times d=50\times 768=0.038$M, and $N\times p\times d=0.46$M parameters, amounting to only 0.04$\%$ and 0.53$\%$ of all ViT-Base parameters, respectively.

\section{Experiments}\label{sec:exp}

We evaluate VPT for a wide range of downstream recognition tasks with pre-trained Transformer backbones across scales. We first describe our experimental setup in~\ref{subsec:evalsetup}, including the pre-trained backbone and downstream tasks, and a brief introduction of alternative transfer learning methods. Then we demonstrate the effectiveness and practical utility of our method in \ref{subsec:exp_results}.
We also systematically study how different design choices would affect performance (\ref{subsec:ablate}),
which leads to an improved understanding of our approach.

\subsection{Experiment Setup}
\label{subsec:evalsetup}

\textbf{Pre-trained Backbones.} We experiment with two Transformer architectures in vision, Vision Transformers (ViT)~\cite{dosovitskiy2020vit} and Swin Transformers (Swin~\cite{liu2021swin}). All backbones in this section are pre-trained on ImageNet-21k~\cite{imagenet_cvpr09}. We follow the original configurations, e.g. , number of image patches divided, existence of \texttt{[CLS]}, etc. . More details are included in~\ref{supsec:detail}.

\textbf{Baselines.}
We compare both variants of VPT with other commonly used fine-tuning protocols:

\begin{enumerate}[nosep, label=(\alph*), font=\small\ttbf,] 
\item \textsc{Full}: fully update \emph{all} backbone and classification head parameters.

\item Methods that focus on the classification head. They treat the pre-trained backbone as a feature extractor, whose weights are fixed during tuning:
\begin{itemize}[leftmargin=0.0em, topsep=0.15mm]
    \item \textsc{Linear}: only use a linear layer as the classification head.
    \item \textsc{Partial}-$k$: fine-tune the last $k$ layers of backbone while freezing the others, as adopted in~\cite{yosinski2014transferable,zhang2016colorful,noroozi2016unsupervised,he2021mae}. It redefines the boundary of backbone and classification head.
    \item \textsc{Mlp}-$k$: utilize a  multilayer perceptron (MLP) with $k$ layers, instead of a linear layer, as classification head.
\end{itemize}

\item Methods that update a subset backbone parameters or add new trainable parameters to backbone during fine-tuning:
\begin{itemize}[leftmargin=0.0em, topsep=0.15mm]
\item \textsc{Sidetune}~\cite{zhang2020side}: train a ``side'' network and linear interpolate between pre-trained features and side-tuned features before being fed into the head.
\item \textsc{Bias}~\cite{cai2020tinytl,zaken2021bitfit}: fine-tune only the bias terms of a pre-trained backbone.
\item \textsc{Adapter}~\cite{houlsby2019parameter,pfeiffer2020adapterfusion,pfeiffer2020AdapterHub}: insert new MLP modules with residual connection inside Transformer layers. 
\end{itemize}
\end{enumerate}

\textbf{Downstream Tasks.} We experiment on the following two collections of datasets:

\textit{FGVC} consists of 5 benchmarked Fine-Grained Visual Classification tasks including CUB-200-2011~\cite{WahCUB_200_2011}, NABirds~\cite{van2015nabirds}, Oxford Flowers~\cite{nilsback2008automated}, Stanford Dogs~\cite{Khosla_FGVC2011dogs} and Stanford Cars~\cite{gebru2017cars}. If a certain dataset only has {\small\texttt{train}} and {\small\texttt{test}} sets publicly available, we randomly split the training set into {\small\texttt{train}} (90\%) and {\small\texttt{val}} (10\%), and rely on {\small\texttt{val}} to select hyperparameters.

\textit{VTAB-1k}~\cite{zhai2019vtab} is a collection of 19 diverse visual classification tasks, which are organized into three groups: 
\textit{Natural} - tasks that contain natural images captured using standard cameras;
\textit{Specialized} - tasks that contain images captured via specialized equipment, such as medical and satellite imagery; and \textit{Structured} - tasks that require geometric comprehension like object counting. Each task of VTAB contains 1000 training examples. Following~\cite{zhai2019vtab}, we use the provided 800-200 split of the {\small\texttt{train}} set to determine hyperparameters and run the final evaluation using the full training data. We report the average accuracy score on {\small\texttt{test}} set within three runs.

We report the average accuracy on the FGVC datasets, and the average accuracy on each of the three groups in VTAB. The individual results on each task are in~\ref{subsec:results_supp}, as are image examples of these aforementioned tasks.

\subsection{Main Results}
\label{subsec:exp_results}

\ref{table:main_vitb} presents the results of fine-tuning a pre-trained ViT-B/16 on averaged across 4 diverse downstream task groups, comparing VPT to the other 7 tuning protocols. We can see that:

\begin{enumerate}[nosep, leftmargin=5mm]
\item \textbf{VPT-Deep outperforms \textsc{Full} (\ref{table:main_vitb}\texttt{(a)}) on 3 out of the 4 problem classes} (20 out of 24 tasks), while using significantly fewer total model parameters (1.18$\times$~\emph{vs}\onedot~24.02$\times$).
Thus, \emph{even if storage is not a concern},  \textsc{VPT} is a promising approach for adapting larger Transformers in vision. Note that this result is in contrast to comparable studies in NLP, where prompt tuning matches, but \emph{does not exceed} full fine-tuning~\cite{lester-etal-2021-power}. 

\item \textbf{VPT-Deep outperforms all the other parameter-efficient tuning protocols (\ref{table:main_vitb}\texttt{(b},\texttt{c)}) across all task groups}, indicating that \textsc{VPT}-\textsc{deep} is the best fine-tuning strategy in storage-constrained environments.

\item Although sub-optimal than \textsc{VPT}-\textsc{deep}, \textsc{VPT}-\textsc{shallow} still offers non-trivial performance gain than head-oriented tuning methods in \ref{table:main_vitb}\texttt{(b)}, indicating that \textsc{VPT}-\textsc{shallow} is a worthwhile choice in deploying multi-task fine-tuned models if the storage constraint is severe.
\end{enumerate}

\subsubsection{VPT on different downstream data size.}
\label{para:exp_datasize}
We look at the impact of training data size on accuracy in the FGVC tasks (VTAB has only 1k training examples). We vary the training data between 10\% and 80\% and compare all methods. The same pre-trained ViT-B is used for downstream training. Task-averaged results for each method on different training data scales are presented in \ref{fig:main_fgvcsize}.

\ref{fig:main_fgvcsize} shows that \textsc{VPT}-\textsc{deep} outperforms all the other baselines across data scales. Digging deeper, methods that use less trainable parameters, i.e. , \textsc{VPT}, \textsc{Linear}, \textsc{Adapter}, \textsc{Bias}, dominate over \textsc{Full} in the low-data regimes. This trend, however, is \emph{reversed} when more training data is available for \textsc{Linear} and \textsc{Adapter}. 
In contrast, \textsc{VPT}-\textsc{deep} still consistently outperforms \textsc{Full} across training data sizes. Although \textsc{Bias} offers similar advantages, it still marginally under-performs \textsc{VPT}-\textsc{deep} across the board (\ref{fig:main_fgvcsize} right).

\subsubsection{VPT on different backbone scales.}
\ref{fig:modelsize} shows VTAB-1k performance under 3 different backbone scales: ViT-\textbf{B}ase/\textbf{L}arge/\textbf{H}uge.
\textsc{VPT}-\textsc{deep} is significantly better than \textsc{Linear} and \textsc{VPT}-\textsc{shallow} across all 3 backbone choices and 3 subgroups of VTAB-1k. More importantly, the advantages of \textsc{VPT}-\textsc{deep} over \textsc{Full} indeed still hold as the model scale increases, i.e. , \textsc{VPT}-\textsc{deep} significantly outperforms \textsc{Full} on \textit{Natural} and \textit{Structured} groups, while offering nearly equivalent performance on \textit{Specialized}.

\subsubsection{VPT on hierarchical Transformers.} We extend \textsc{VPT} to Swin~\cite{liu2021swin}, which employs MSA within local shifted windows and merges patch embeddings at deeper layers. For simplicity and without loss of generality, we implement VPT in the most straightforward manner: the prompts are attended within the local windows, but are ignored during patch merging stages. The experiments are conducted on the~ImageNet-21k supervised pre-trained Swin-\textbf{B}ase. 
\textsc{VPT} continues to outperform other parameter-efficient fine-tuning methods (\texttt{b, c}) for all three subgroups of VTAB \ref{table:main_swinb}, though in this case \textsc{Full} yields the highest accuracy scores overall (at a heavy cost in total parameters). 

It is surprising that the advantage of \textsc{VPT}-\textsc{deep} over \textsc{VPT}-\textsc{shallow} diminishes for \textit{Natural}: \textsc{VPT}-\textsc{shallow} yields slightly better accuracy scores than full fine-tuning.

\subsection{Ablation on Model Design Variants}
\label{subsec:ablate}
We ablate different model design choices on the supervised ImageNet-21k pre-trained ViT-Base and evaluate them on VTAB, with same setup in \ref{table:main_vitb}. See more in~\ref{supp_analysis}.

\subsubsection{Prompt Location.} An important distinction between \textsc{VPT} and other methods is the extra learnable parameters introduced as \textit{inputs} for the Transformer layers. 
\ref{fig:ablate_loc} ablates different choices on how and where to insert prompts in the input space, and how they would affect the final performance.

\textit{Prepend or Add?} Instead of prepending prompts to the sequence of the image patches embeddings $\vec{E}_i$ as described in~\ref{subsec:method_vp}, 
another option is to directly \emph{add} prompts element-wise to those embeddings, keeping the Transformer's input sequence length the same as before. Though this variant is competitive to \textsc{Full} in some cases (e.g. , VTAB-\textit{Natural}), its performance generally falls behind with the default \texttt{Prepend} in both deep and shallow settings. More discussion on this phenomenon is in~\ref{app:seq}.

\textit{Latent or pixel space?} Instead of inserting the prompts as latent vectors for the first Transformer layer, one could introduce prompts in the \textit{pixel} level before the \texttt{Embed} layer in~\ref{eq:embed},
i.e. , \texttt{Prepend-pixel} and \texttt{Concat-channel}.
\ref{fig:ablate_loc} shows that the adaption performance \emph{decreases} for these two variants. For example, the accuracy score of prepending shallow prompts before the projection layer (\texttt{Prepend-pixel}) drops 6.9$\%$, compared to the default prepending in the embedding space (\texttt{Prepend}) on VTAB-\textit{Natural}. The performance further deteriorates (even as large as 30 accuracy scores drop on VTAB-\textit{Natural}) if we instead concatenate a new channel to the input image (\texttt{Concat-channel}). These observations suggest that it's easier for prompts to learn condensed task-dependent signals in the latent input space of Transformers.

\subsubsection{Prompt Length.} This is the only additional hyper-parameter needed to tune for VPT compared to full fine-tuning. For easy reference, we also ablate two other baselines on their individual additional hyper-parameters, i.e. , number of layers for \textsc{Mlp} and reduction rate for \textsc{Adapter}. As shown in \ref{fig:ablate_length}, the optimal prompt length varies across tasks. Notably, even with as few as only \emph{one} prompt, \textsc{VPT}-\textsc{deep} still significantly outperforms the other 2 baselines, and remains competitive or even better compared to full fine-tuning on VTAB-\textit{Structured} and \textit{Natural}.

\subsubsection{Prompt Depth.}
\ref{fig:ablate_depth} ablates which and how many layers to insert prompts. Each variant reports the best prompt length selected with {\small\texttt{val}} set.
\textsc{VPT}'s performance is positively correlated with the prompt depth in general. Yet the accuracy drops if we insert prompts from \textcolor{prompt_red}{top to bottom}, suggesting that prompts at earlier Transformer layers matter more than those at later layers. 

\subsubsection{Final Output.} Following the original configuration of ViT, we use the final embedding of $\texttt{[CLS]}$, i.e. , $\vec{x}_N$, as the classification head input, which is also the default setting in our ViT experiments. As shown in \ref{fig:ablate_output}, 
if we use the average pooling on image patch output embeddings 
$\vec{E}_N$ as final output (\texttt{Image-pool}), the results essentially remain the same (e.g. , 82.4 \emph{vs}\onedot~82.3 for VTAB-\textit{Specialized}).
However, if the pooling involves final prompt outputs $\vec{Z}_N$ (\texttt{Prompt-pool} and \texttt{Global-pool}), the accuracy could drop as large as 8 points.

\section{Analysis and Discussion}
\label{sec:ana}

\subsubsection{Visualization.}
\ref{fig:tsne} shows t-SNE~\cite{van2008visualizing} visualizations of $\mathbf{x_N}$, i.e. , embeddings of \texttt{[CLS]} after the last Transformer layer and before the classification head, for 3 tasks in VTAB (SVNH~\cite{netzer2011reading}, EuroSAT~\cite{helber2019eurosat}, Clevr/count~\cite{johnson2017clevr}), one for each subgroup. All plots show that \textsc{VPT}-\textsc{deep} enables linearly separable representations while using less parameters than \textsc{Full}.
We also observe that extra tunable parameters for every Transformer layer (\textsc{VPT}-\textsc{deep}) improve the performance, compared to \textsc{VPT}-\textsc{shallow}, which only inserts prompts for the first layer's input. Interestingly on Clevr/count (\ref{fig:tsne}(c)), \textsc{VPT}-\textsc{deep} and \textsc{Full} recover the underlying manifold structure of the task (counting objects in images~\emph{vs}\onedot street number or landscape recognition), unlike \textsc{VPT}-\textsc{shallow} and \textsc{Linear}.

\subsubsection{Apply VPT to more vision tasks.} We explore the feasibility of \textsc{VPT} beyond visual classification, by evaluating ADE20K~\cite{zhou2019ade20} semantic segmentation task with a Transformer model, SETR-PUP~\cite{zheng2020rethinking}. It adds a standard ConvNet head to the ViT backbone to perform segmentation. The de-facto approach is still fully fine-tuning the pre-trained backbone together with the ConvNet head (\textsc{Full}). We include two more protocols for comparison: only update the head layers (\textsc{Head Only}), update head layers and bias vectors in the backbone (\textsc{Bias}). 
In~\ref{table:main_seg}, we report {\small\texttt{val}} mIoU results with and without multi-scale inference. Though parameter-efficient protocols could not compete with \textsc{Full}, \textsc{VPT} is still comparable with \textsc{Bias}.
Notably, \textsc{VPT} offers competitive results to a fully fine-tuned state-of-the-art ConvNet model (DeepLab v3+~\cite{chen2018encoder}), while tuning significantly less parameters (15M \emph{vs}\onedot 64M, respectively).

\subsubsection{Apply VPT to more pre-training methods.} In addition to the backbones pre-trained with labeled data, we experiment with two self-supervised objectives: MAE~\cite{he2021mae} and MoCo v3~\cite{chen2021mocov3}.
\ref{table:main_ssl} reports the results on VTAB-1k with ViT-B. We observe that both variants of \textsc{VPT} surpass \textsc{Linear}, yet the comparisons among other techniques are less conclusive. For MAE, other parameter-efficient methods, e.g. , \textsc{Partial}-1, outperform both \textsc{VPT} and \textsc{Linear}.
In the case of MoCo v3, \textsc{VPT} no longer holds the best performance, though it is still competitive with the others. This suggests that these two self-supervised ViTs are fundamentally different from the supervised ones in previous sections. Exactly why and how these differences arise remain open questions.

\subsubsection{Apply VPT to ConvNets.} We examine the idea of adding trainable parameters in the input space of ConvNets: padding both height and width by $p$ learnable prompt pixels for the input image. Though this operation seems unconventional, we implement VPT this way given there is no obvious solution to add location-invariant prompts similar to the Transformer counterparts. In fact this approach has been explored before in the adversarial attack literature~\cite{elsayed2018adversarial}. The value of $p$ in our experiment is 2 orders of magnitude smaller than previous work: e.g. , 5 \emph{vs}\onedot 263. Most importantly, we cast this idea in the lens of transfer learning. See~\ref{vpt_ar} for more discussion.

\ref{table:main_rnxrn} presents the results for ConvNeXt-B~\cite{liu2022convnet} (pre-trained on ImageNet-21k) and ResNet-50~\cite{he2016rn} (pre-trained on ImageNet-1k), respectively.
\textsc{VPT} works well in a larger ConvNet backbone, ConvNeXt-B, offering accuracy gains over other sparse tuning protocols (\texttt{b}, \texttt{c}), and outperforming \textsc{Full} on 8 out of 19 cases. The advantages of \textsc{VPT}, however, diminish with smaller ConvNet (ResNet-50), as there is no clear winner for all 19 VTAB-1k tasks.

\section{Conclusion}

We present Visual Prompt Tuning, a new parameter-efficient approach to leverage large vision Transformer models for a wide range of downstream tasks. VPT introduces task-specific learnable prompts in the input space, keeping the pre-trained backbone fixed. We show that VPT can surpass other fine-tuning protocols (often including full fine-tuning) while dramatically reducing the storage cost. Our experiments also raise intriguing questions on fine-tuning dynamics of vision Transformers with different pre-training objectives, and how to transfer to broader vision recognition tasks in an efficient manner. We therefore hope our work will inspire future research on how best to tap the potential of large foundation models in vision.

\small{
\textbf{Acknowledgement.}
Menglin is supported by a Meta AI research grant awarded to Cornell University, Luming and Bharath is supported by NSF IIS-2144117, Serge is supported in part by the Pioneer Centre for AI, DNRF grant number P1. We would like to thank Alexander Rush, Yin Cui for valuable suggestions and discussion. }

\end{document}
